\documentclass[12pt, a4paper]{article}

\usepackage[margin=1in]{geometry}
\usepackage{amsmath, amssymb}
\usepackage{graphicx}
\usepackage{booktabs}
\usepackage{hyperref}
\usepackage{parskip}
\usepackage{xcolor}

\title{\textbf{Demo \LaTeX{} Document} \\ \large A Quick Tour of Common Features}
\author{Daanish Muzaffar}
\date{\today}

\begin{document}

\maketitle
\tableofcontents
\newpage

% ─────────────────────────────────────────────
\section{Introduction}

This document demonstrates common \LaTeX{} features: text formatting,
mathematics, tables, and lists. \LaTeX{} separates \textit{content} from
\textit{presentation} --- you write markup and the compiler handles typesetting.

\textbf{Bold}, \textit{italic}, \texttt{monospace}, and
\textcolor{blue}{coloured text} are all trivial.

% ─────────────────────────────────────────────
\section{Mathematics}

Inline math sits in the flow of text: $E = mc^2$, or the quadratic formula
$x = \frac{-b \pm \sqrt{b^2 - 4ac}}{2a}$.

Display equations get their own line and are automatically numbered:

\begin{equation}
    \int_{-\infty}^{\infty} e^{-x^2}\, dx = \sqrt{\pi}
\end{equation}

\begin{equation}
    \hat{\sigma}^2 = \frac{1}{T} \sum_{t=1}^{T} \min(r_t - \text{MAR},\; 0)^2
\end{equation}

Aligned multi-line equations use the \texttt{align} environment:

\begin{align}
    \text{Sharpe}  &= \frac{\mu_r - r_f}{\sigma_r} \cdot \sqrt{252} \\[4pt]
    \text{Sortino} &= \frac{\mu_r - \text{MAR}}{\hat{\sigma}} \cdot \sqrt{252}
\end{align}

% ─────────────────────────────────────────────
\section{Lists}

\subsection{Bullet list}
\begin{itemize}
    \item MA Crossover --- trend-following via short/long moving average gap
    \item RSI --- mean-reversion triggered by oversold/overbought thresholds
    \item Donchian Channel --- breakout when price exceeds the $N$-day high
\end{itemize}

\subsection{Numbered list}
\begin{enumerate}
    \item Download price data
    \item Split into in-sample and out-of-sample windows
    \item Optimise parameters on in-sample Sortino ratio
    \item Evaluate frozen parameters on out-of-sample data
\end{enumerate}

% ─────────────────────────────────────────────
\section{Tables}

\begin{table}[h]
    \centering
    \caption{Signal-to-Basket Assignment}
    \begin{tabular}{llll}
        \toprule
        \textbf{Signal} & \textbf{Basket} & \textbf{ETF} & \textbf{IS Sortino} \\
        \midrule
        MA Crossover      & Financials   & XLF & 1.11 \\
        RSI               & Technology   & XLK & 0.94 \\
        Donchian Channel  & Health Care  & XLV & 0.87 \\
        \bottomrule
    \end{tabular}
    \label{tab:signals}
\end{table}

% ─────────────────────────────────────────────
\section{Useful Environments}

A block quote for referencing sources:

\begin{quote}
    ``The Sortino ratio penalises only downside deviation, making it more
    appropriate than the Sharpe ratio for asymmetric return distributions.''
    \hfill --- Sortino \& van der Meer (1991)
\end{quote}

Verbatim code:

\begin{verbatim}
def donchian_signal(prices, window=55):
    high = prices.rolling(window).max().shift(1)
    return (prices > high).astype(int)
\end{verbatim}

% ─────────────────────────────────────────────
\section{Conclusion}

That covers the essentials. Compile with:
\begin{center}
    \texttt{pdflatex demo.tex}
\end{center}
or use an online editor like \href{https://www.overleaf.com}{Overleaf} if you
don't have a local \LaTeX{} distribution installed.

\end{document}
